\documentclass[11pt]{article}
\usepackage[utf8]{inputenc}
\usepackage[T1]{fontenc}
\usepackage{amsmath,amssymb,amsfonts}
\usepackage{tgtermes}
\usepackage{newtxmath}
\usepackage[version=4]{mhchem}
\usepackage{siunitx}
\usepackage[margin=1in]{geometry}
\emergencystretch=2em
\DeclareSIUnit{\Molar}{M}
\newcommand{\epss}{\varepsilon_{\mathrm{s}}}
\newcommand{\Rtheta}{R_{\theta}}
\newcommand{\Tpore}{T_{\mathrm{pore}}}
\begin{document}
\section*{Figure captions}
% AUTO-GENERATED FROM CURRENT SOURCES. main_R582.tex SHA-256 8AE9FF4F244E9D5AB2308D88E3B1E98BC5C9E946931F2CD7EF6A9CAA7D3E606D; SI_R582.tex SHA-256 8BDC18C71088398E0E9B1A55359FA3FDF688DE536B49C2C6C497EA34C47F6C1C.
% Caption bodies are verbatim source extractions. Regenerate after every source-caption edit.
\section*{Main manuscript}

\noindent\textbf{Figure 1.}\begingroup\space
% BEGIN_CAPTION source=main number=1 label=fig:experimental-problem
\textbf{Full-cell records motivate a positive-electrode state analysis.}
\textbf{a,} The measured cell response contains both electrodes, whereas the model and interpretation focus on the porous positive carbon felt.
\textbf{b,} Representative cycle-20 charge at 40 mA cm$^{-2}$. The thin grey trace is the registered interpolated voltage, the dark trace is its frozen cubic Savitzky--Golay smooth, and the lower trace is $\mathrm{d}V/\mathrm{d}Q$. The shaded interval contains the voltage-rise feature. $Q_s=83.0$ and $Q_{f,\mathrm{cal}}=99.6$ mAh cm$^{-2}$ are pre-existing calibrated-model coordinates; neither was fitted to the derivative landmarks, and their overlay is not independent validation.
\textbf{c,} Coulombic efficiency for the same-cell pristine 80--360 mA cm$^{-2}$ rate ladder. Symbols are within-cell medians, and vertical bars span the full within-cell minimum--maximum range across available sequential cycles. The separate 40 mA cm$^{-2}$ pristine cell and the 400 mA cm$^{-2}$ P350C-positive/pristine-negative proxy are deliberately unconnected.
\textbf{d,} Primary-window rise onset and maximum for cycles 20--23 from one physical cell. Pale spans give extrema across three prespecified derivative windows and are estimator-sensitivity ranges, not confidence intervals. These records motivate, but do not independently validate, the positive-electrode state model.
% END_CAPTION source=main number=1 label=fig:experimental-problem
\par\endgroup\medskip

\noindent\textbf{Figure 2.}\begingroup\space
% BEGIN_CAPTION source=main number=2 label=fig:domain-state
\textbf{A state-resolved model of the ZIFB positive electrode.}
\textbf{a,} The two-dimensional carbon-felt domain spans $x=3$--5 mm from the current collector to the separator and $y=0$--20 mm along electrolyte flow.
\textbf{b,} Free-\ce{I2} supersaturation, $S$, gates retained solid, $\varepsilon_s$. The calibrated relation gives $R_\theta=1-\theta_{\mathrm{cal}}$; the native reaction-area multiplier is $A_{\mathrm{bare}}/A_0=R_\theta T_{\mathrm{pore}}$, and $V$ is a model output.
\textbf{c,} The baseline simulation reaches average saturation at $Q_s=83.0$ mAh cm$^{-2}$ and 50\% calibrated accessibility loss at $Q_{f,\mathrm{cal}}=99.6$ mAh cm$^{-2}$. The state annotations do not define local onset or microscopic deposit form.
% END_CAPTION source=main number=2 label=fig:domain-state
\par\endgroup\medskip

\noindent\textbf{Figure 3.}\begingroup\space
% BEGIN_CAPTION source=main number=3 label=fig:spatial-progression
\textbf{Retained solid and native remaining area redistribute reaction within the positive electrode.}
Continuum-model snapshots are shown at $Q=80$, 100 and 120 mAh cm$^{-2}$, respectively before the electrode-average saturation point, near the calibrated half-accessibility point and at late charge. The capacity rail distinguishes the exported snapshots from $Q_s=83.0202$ and $Q_{f,\mathrm{cal}}=99.5901$ mAh cm$^{-2}$; $Q=100$ mAh cm$^{-2}$ is close to, but not identical to, $Q_{f,\mathrm{cal}}$. The map rows show retained-solid fraction, $\varepsilon_s$, native remaining bare-area fraction, $A_{\mathrm{bare}}/A_0$, and total reaction-current density, $j_{\mathrm{total}}$, using one scale per quantity across all capacities. The lower plot gives the arithmetic mean of the native remaining-area field over the uniformly sampled flow coordinate $y$. Six negative-current nodes at $Q=100$ mAh cm$^{-2}$, spanning $-5.72\times10^{-4}$ to $-1.80\times10^{-4}$ A m$^{-2}$, are retained as a neutral mask rather than clipped or converted to magnitudes. Every map uses the same 3--5 mm through-plane and 0--20 mm flow domain. These are continuum-model fields, not micrographs or evidence for a unique deposit morphology or pore-blocking front.
% END_CAPTION source=main number=3 label=fig:spatial-progression
\par\endgroup\medskip

\noindent\textbf{Figure 4.}\begingroup\space
% BEGIN_CAPTION source=main number=4 label=fig:closure-identifiability
\textbf{Accessibility relation controls the late charge response.}
\textbf{a,} Calibration-controlled remaining-accessibility factor, $R_\theta=1-\theta$, for the calibrated baseline, a one-way island-model shadow evaluated on the baseline solid inventory, and the fully coupled island-model variant. The matched simulations give $Q_s=82.890$ and $82.892$ mAh cm$^{-2}$.
\textbf{b,} Voltage trajectories and the island-minus-baseline difference, which reaches $-288.2$ mV at 120 mAh cm$^{-2}$.
\textbf{c,} Each solved branch's solid-\ce{I2} fraction; horizontal lines mark the solid-fraction values at $R_\theta=0.5$ for the baseline and island relation.
\textbf{d,} Accessibility--coefficient pairs that reproduce the same selected voltage contribution. This controlled model-form sensitivity does not identify deposit morphology.
% END_CAPTION source=main number=4 label=fig:closure-identifiability
\par\endgroup\medskip

\noindent\textbf{Figure 5.}\begingroup\space
% BEGIN_CAPTION source=main number=5 label=fig:multiscale-bounds
\textbf{Independent bounds on positive-electrode model inputs.}
\textbf{a,} Periodic CP2K BSSE-corrected adsorption energies for dry single-\ce{I2} placement, reported relative to basal carbon.
\textbf{b,} Five-SOC bulk-MD means for \ce{I^-}, \ce{I3^-} and \ce{I2Br^-} under $q=0.8$ ECC and formal-charge parameterizations in the representative \SI{1}{M}~\ce{ZnI2}$+$\SI{4}{M}~\ce{NH4Br} supporting electrolyte; error bars report propagated within-trajectory stability. Shading denotes the formal-charge low-end carrier-proxy range ($0.425$--$0.664\times10^{-9}~\si{m^2.s^{-1}}$), and the dashed line marks the continuum baseline $D_{\mathrm{eff}}=0.50\times10^{-9}~\si{m^2.s^{-1}}$.
\textbf{c,} Geometric remaining area $A/A_0$ for sparse ($N=10^{11}~\si{m^{-2}}$) and dense ($N=10^{14}~\si{m^{-2}}$) single-fibre placement; lines connect computed nodes, including the exact zero-solid boundary.
\textbf{d,} Six-law permeability envelope for an idealized pore network over the baseline solid-\ce{I2} range; the inset retains the full computed range. At $\varepsilon_s=3.219\times10^{-3}$, $K/K_0=0.981$--$0.992$. These lower-scale calculations bound transport and accessibility separately. They do not determine the voltage-calibrated inventory-to-accessibility relation or identify deposit morphology or pore closure.
% END_CAPTION source=main number=5 label=fig:multiscale-bounds
\par\endgroup\medskip

\noindent\textbf{Figure 6.}\begingroup\space
% BEGIN_CAPTION source=main number=6 label=fig:operating-levers
\textbf{Applied current and oxidized-carrier diffusivity shift the averaged saturation point.}
\textbf{a,} Positive-electrode-average $S=1$ capacity, $Q_s$, versus applied current. Filled circles denote exact crossings from full continuum simulations. The line connects only the complete 20 and 40 mA cm$^{-2}$ traces. The open triangle at 80 mA cm$^{-2}$ denotes the right-censored result $Q_s>40$ mAh cm$^{-2}$; the separate 120 mA cm$^{-2}$ calculation crosses at 8.39 mAh cm$^{-2}$.
\textbf{b,} $Q_s$ across seven full $D_{\mathrm{eff}}/D_0$ simulations; adjacent points are joined without fitting or extrapolation. Shading marks the MD-informed 0.8--1.4 $D_0$ range.
\textbf{c,} Signed shifts from the baseline $Q_s=83.0202$ mAh cm$^{-2}$. Circles, dashed diamonds and open dotted triangles distinguish full continuum simulations, an analytical boundary-layer scenario and an MD-informed range mapped through the solved diffusivity response. The analytical flow row uses $\delta\propto v^{-0.5}$, and the felt comparison changes thickness and porosity together. The smooth-permeability replacement is reported separately through its $+1.27$ mV endpoint-voltage difference.
% END_CAPTION source=main number=6 label=fig:operating-levers
\par\endgroup\medskip

\section*{Supplementary Information}

\noindent\textbf{Supplementary Figure S1.}\begingroup\space
% BEGIN_CAPTION source=si number=1 label=fig:si-derivative
\textbf{Derivative-window sensitivity within one physical cell.} \textbf{a--d,} Full-cell charge-
curve derivatives for eligible cycles 20--23 of one pristine cell. Direct labels identify the three
prespecified physical windows; open circles mark registered rise onsets and filled triangles the registered
maxima. Boundary fluctuations outside the maximum-search interval remain visible. All panels share one
derivative scale. The cycles are within-cell repeated observations, not independent replicates. These
full-cell features do not identify a positive-electrode internal state. Derived electrolyte labels follow
\texttt{EXP-META-001}; raw records are unchanged.
% END_CAPTION source=si number=1 label=fig:si-derivative
\par\endgroup\medskip

\noindent\textbf{Supplementary Figure S2.}\begingroup\space
% BEGIN_CAPTION source=si number=2 label=fig:si-composition
\textbf{Cell-level descriptive capacity envelope across the existing composition records.} Each
filled marker is one physical cell and gives the maximum, across commanded levels, of the within-level median
delivered capacity. Vertical segments are full sequential-cycle ranges at the selected command, and short
black bars are condition medians. Physical-cell counts are 1/1/1/2/1 at 0/1/2/3/4~M NH4Br. The two
acquisitions at 2~M form one conservatively stitched cell; the two 3~M records are distinct, and one 3~M
endpoint is right-censored. The display is descriptive and does not identify an internal electrode state.
% END_CAPTION source=si number=2 label=fig:si-composition
\par\endgroup\medskip

\noindent\textbf{Supplementary Figure S3.}\begingroup\space
% BEGIN_CAPTION source=si number=3 label=fig:si-compression
\textbf{Protocol-bounded capacity brackets across existing felt-thickness cells.} \textbf{a,} Primary
85\% within-cell median-CE endpoint. The filled blue points connect only the three explicit March cells; the
open 2.0~mm cross-batch proxy is unconnected. \textbf{b,} Prespecified 80\%, 85\% and 90\% endpoint
sensitivity. Each bracket is program resolution, not a confidence interval. The data do not define a
population scaling law, pore-volume law or morphology relation.
% END_CAPTION source=si number=3 label=fig:si-compression
\par\endgroup\medskip

\noindent\textbf{Supplementary Figure S4.}\begingroup\space
% BEGIN_CAPTION source=si number=4 label=fig:si-state-fields
\textbf{Six-capacity inventory of modeled positive-electrode state and reaction fields.} Columns show
$Q=0$, 80, 96, 100, 110 and \SI{120}{mAh.cm^{-2}}. \textbf{a,} Free-\ce{I2} stress $S$ on a shared scale
centred at $S=1$. \textbf{b,} Native retained-solid fraction $\epss$ on a shared symlogarithmic percentage
scale. \textbf{c,} Native remaining reaction-area fraction $A_{\mathrm{bare}}/A_0=\Rtheta\Tpore$ from the
registered \texttt{A\_bare\_frac} export. \textbf{d,} Total reaction-current density. Positive values use one
logarithmic scale; non-positive nodes are neutrally masked without taking absolute values. Their counts are
248, 0, 33, 6, 0 and 0 in column order, and signed values remain in the source table. Each map contains 5995
nodes and uses nearest-neighbour rendering without smoothing or clipping. These fields are model outputs, not
microscopy, deposit morphology, a blockage front or electrical potential.
% END_CAPTION source=si number=4 label=fig:si-state-fields
\par\endgroup\medskip

\noindent\textbf{Supplementary Figure S5.}\begingroup\space
% BEGIN_CAPTION source=si number=5 label=fig:si-hydraulic-fields
\textbf{Hydraulic response across the same six positive-electrode states.} \textbf{a,} Smooth
relative permeability $K/K_0$. \textbf{b,} Native velocity magnitude in \si{m.s^{-1}}. \textbf{c,} Hydraulic
pressure in Pa relative to $p_{\mathrm{out}}=0$. Each row has one fixed scale across all capacities, and every
map contains the same 5995 exported nodes. The pressure field is hydraulic, not an electrical potential. The
smooth permeability channel does not resolve local pore closure or blockage.
% END_CAPTION source=si number=5 label=fig:si-hydraulic-fields
\par\endgroup\medskip

\noindent\textbf{Supplementary Figure S6.}\begingroup\space
% BEGIN_CAPTION source=si number=6 label=fig:si-voltage-degeneracy
\textbf{Full-cell voltage fit and accessibility--coefficient degeneracy.} \textbf{a,} Selected
pristine full-cell charge trace and reduced observation-layer fit. The shaded interval contains 72 points from
2.672 to \SI{97.540}{mAh.cm^{-2}} in
$V=c_0+c_1Q+b_{\mathrm{eff}}\ln(1/\Rtheta)$; the lower axis shows fit minus data. The fit gives
$c_0=\SI{1.33554}{V}$, $c_1=9.819\times10^{-4}~\si{V.(mAh.cm^{-2})^{-1}}$,
$b_{\mathrm{eff}}=\SI{0.13622}{V}$ and RMSE \SI{7.24}{mV}. \textbf{b,} End-point $\Rtheta$--coefficient
pairs that preserve the selected \SI{70.813}{mV} accessibility-related voltage contribution. Open circles
are registered pairs and the line is the analytical relation. Here $\Rtheta$ is a reduced observation-layer
factor, not native remaining reaction area. Voltage therefore does not identify the inventory-to-accessibility
relation or deposit morphology uniquely.
% END_CAPTION source=si number=6 label=fig:si-voltage-degeneracy
\par\endgroup\medskip

\noindent\textbf{Supplementary Figure S7.}\begingroup\space
% BEGIN_CAPTION source=si number=7 label=fig:si-single-i2
\textbf{Relative single-\ce{I2} BSSE-corrected adsorption-energy ordering and optimized
geometries.} \textbf{a,} BSSE-corrected adsorption energies relative to the basal calculation.
\textbf{b--e,} Exact optimized
geometries in the same site order; displayed adjacency guides are distance based and do not assign bond order.
Within the declared finite periodic calculation, the tested C--OH and C=O motifs lie below basal carbon and
the vacancy lies slightly above it. The unavailable raw single-\ce{I2} charge-density fields are not used.
% END_CAPTION source=si number=7 label=fig:si-single-i2
\par\endgroup\medskip

\noindent\textbf{Supplementary Figure S8.}\begingroup\space
% BEGIN_CAPTION source=si number=8 label=fig:si-two-i2
\textbf{Compact and separated two-\ce{I2} configurations at the tested C--OH site.} \textbf{a,}
Registered compact-minus-separated electronic-energy comparison. \textbf{b,c,} Exact optimized compact and
separated geometries. \textbf{d,} Recoverable two-\ce{I2} charge-density-difference field, thresholded at
$\Delta\rho\geq+0.002$ or $\leq-0.002~e\,\text{\AA}^{-3}$ without smoothing or interpolation; 428
accumulation and 78 depletion points are retained. One periodic translation is used only to display the
compact cluster continuously. The compact configuration is \SI{0.689291}{eV} below the separated reference
within this calculation. The field is a diagnostic, not a reaction coordinate, solution equilibrium or
nucleation pathway.
% END_CAPTION source=si number=8 label=fig:si-two-i2
\par\endgroup\medskip

\noindent\textbf{Supplementary Figure S9.}\begingroup\space
% BEGIN_CAPTION source=si number=9 label=fig:si-md-ladder
\textbf{Full carrier-mobility ladder under charge-scaled and formal-charge parameterizations.}
\textbf{a,} $q=0.8$ ECC. \textbf{b,} $q=1.0$ formal charges. The condition is 1~M \ce{ZnI2} + 4~M NH4Br
at \SI{298.15}{K}. Symbols are means across five different SOC compositions; thin spans show their min--max
range and capped bars the propagated four-block within-trajectory stability measure. Neither is replica
uncertainty. Oxidized-iodine carrier mobilities lie below iodide in both parameterizations, while absolute
values remain force-field dependent.
% END_CAPTION source=si number=9 label=fig:si-md-ladder
\par\endgroup\medskip

\noindent\textbf{Supplementary Figure S10.}\begingroup\space
% BEGIN_CAPTION source=si number=10 label=fig:si-comparator-definitions
\textbf{Flat definitions of the single-fibre and pore-network comparator models.} \textbf{a,}
Orthographic views of the prescribed cylindrical unit cell
($r_f=\SI{7.50}{\micro m}$, $r_{\mathrm{out}}=\SI{19.36}{\micro m}$,
$L_z=\SI{150}{\micro m}$, $\varepsilon=0.85$). Normalized concentration is fixed at one at the outer radius;
a bare wall node has $c=0$, a covered node has zero normal flux, and azimuthal/axial directions are periodic.
\textbf{b,} Exact $z=7$ slice of the registered $15\times15\times15$ cubic network (225 nodes and 420 in-slice
throats). Line width maps initial throat radius; opposing faces have normalized Dirichlet values one and zero.
The \SI{60}{\micro m} lattice and displayed equations define conductance, allocation and radius update. These
are prescribed comparator domains, not images or reconstructions of the felt. No deposit field, solved
pressure, electrical potential, flux or blockage front is shown.
% END_CAPTION source=si number=10 label=fig:si-comparator-definitions
\par\endgroup\medskip

\noindent\textbf{Supplementary Figure S11.}\begingroup\space
% BEGIN_CAPTION source=si number=11 label=fig:si-accessibility-families
\textbf{Remaining-accessibility families across the full tested placement-density range.}
\textbf{a,} Geometric single-fibre $A/A_0$ for all four registered comparator families. Lines connect five
computed nodes without fitting. The dashed grey curve is the separately calibrated continuum $\Rtheta$
relation sampled at the registered inventory nodes; it is not a geometric family and not native
$A_{\mathrm{bare}}/A_0$. \textbf{b,} Five shared inventory nodes for the two $\phi_{\mathrm{ppt}}=0.005$
families. Increasing prescribed placement density from $10^{11}$ to $10^{13}~\si{m^{-2}}$ changes geometric
remaining area by 0.085--0.309 over $\epss=3.39\times10^{-4}$ to $3.39\times10^{-3}$. These deterministic
comparators are not measured coverage, felt morphology or validation of the calibrated relation.
% END_CAPTION source=si number=11 label=fig:si-accessibility-families
\par\endgroup\medskip

\noindent\textbf{Supplementary Figure S12.}\begingroup\space
% BEGIN_CAPTION source=si number=12 label=fig:si-flow-postprocess
\textbf{Conditional flow dependence of the analytical boundary-layer scenario.} The registered
baseline was postprocessed with $\delta/\delta_0=(v/v_0)^{-m}$ for $m=0.4$, 0.5 and 0.6, with
$\delta_0=\SI{25}{\micro m}$ at \SI{50}{mL.min^{-1}}. All curves pass through
$Q_s=\SI{83.020}{mAh.cm^{-2}}$ at nominal flow. The shaded band marks \SIrange{25}{100}{mL.min^{-1}}; the
$m=0.5$ scenario moves $Q_s$ from 74.992 to \SI{88.733}{mAh.cm^{-2}}. These are registered analytical
evaluations that rescale only sub-grid generation stress, not a full flow-dependent boundary-layer solve or a
measured mass-transfer law.
% END_CAPTION source=si number=12 label=fig:si-flow-postprocess
\par\endgroup\medskip

\noindent\textbf{Supplementary Figure S13.}\begingroup\space
% BEGIN_CAPTION source=si number=13 label=fig:si-permeability
\textbf{Voltage response to replacing the smooth permeability relation.} \textbf{a,} Voltage
difference between the re-solve using the pore-network-derived smooth path and the matched Kozeny--Carman
baseline. The maximum absolute difference is \SI{5.214}{mV} at \SI{117.444}{mAh.cm^{-2}}; the difference at
\SI{120}{mAh.cm^{-2}} is $+\SI{1.271}{mV}$. \textbf{b,} Corresponding smooth $K/K_0$ paths; endpoint values
are 0.98250 and 0.95629. This control addresses smooth bulk permeability only. Local pore-throat blockage is
not resolved or excluded.
% END_CAPTION source=si number=13 label=fig:si-permeability
\par\endgroup\medskip

\end{document}
